\section{Introduction}\label{sec:intro}

\subsection{Crushing as an energy and capacity decision}\label{sec:intro-context}

Crushing is the first stage of comminution, and its setting decides how much size reduction is left to the grinding mill
downstream. In Chilean copper mining the concentrator, the plant that holds the crushing and grinding circuits, took
$56{,}970$~TJ of electricity in 2024, $55.7\%$ of the $102{,}307$~TJ the sector used, and the national copper commission
states that this electricity goes especially to the crushing and grinding operations~\cite{cochilco2025}. The electricity
per tonne of fine copper in concentrates rose from $8{,}946$ to $13{,}079$~MJ between 2010 and 2024, an increase of
$46.2\%$, while the national average ore grade fell from $0.75\%$ to $0.62\%$; the same report ties the rise to deeper
deposits and harder rock that directly affect comminution~\cite{cochilco2025}. The projection of the same agency for
2025 to 2034 has the electricity of the copper sector growing by $20.2\%$, from $27.6$ to $33.2$~TWh, against $8.3\%$ for
fine copper, with concentration alone at $18.1$~TWh in 2034~\cite{cochilco2026}. Outside Chile, an analysis of global
metallic mining reports a decline of average copper ore grades of about a quarter over ten years together with a rise in
energy use larger than the rise in production~\cite{calvo2016}, and the International Energy Agency names declining ore
grades among the reasons for an implied supply deficit of copper in 2035 under its stated-policies scenario~\cite{iea2025}.

The energy is not spent efficiently. The United States Department of Energy bandwidth study of mining (data year 2000)
attributes $494$ of $1{,}246$ trillion Btu per year to grinding and $52$ to crushing, and records that one percent is the
most common estimate of the efficiency of grinding against about fifty percent for crushing on a mechanical basis, with
the warning that the estimates depend on the metric~\cite{doe2007}. An exergy analysis of the comminution of a central-Chile
copper ore in a laboratory ball mill finds an efficiency of about three percent against the Hukki and Morrell
benchmark~\cite{jacob2017}. A coarser crusher product passes work to grinding, the more energy-intensive step, so the
setting of a crusher is a decision about the energy of the whole circuit and not only about the crusher.

\subsection{Who decides, and how it is decided today}\label{sec:intro-roles}

Three roles carry the decision. The comminution or process engineer chooses the flowsheet, the machines and their
settings at design and at each review; the plant operator and the control-room supervisor move the setting and the speed
day to day; the maintenance and reliability team decides when the liners are changed. Four families of tools serve them.
Empirical energy laws, of which the Bond work index is the best known, size machines and estimate energy and have been
used for about sixty years in mill design, but they cannot give a full circuit mass balance or a size
distribution~\cite{lacerda2016}, and design test data and benchmarks are often modified, worn or misapplied in
practice~\cite{bailey2009}. Population-balance crusher models in the Whiten and Andersen tradition fit industrial
surveys closely when calibrated: a review classifies $61$ works on cone-crusher models and control~\cite{yamashita2021},
a calibration with industrial data from two iron-ore mines treats every entry of the breakage matrix as a fitted
parameter~\cite{duarte2021}, and a fit-for-purpose model of a Metso HP500 secondary cone was calibrated on ten surveys
over several settings and feeds~\cite{rocha2024}. Discrete-element simulation reproduces the coarse part of the product
and the control signals of a crusher with relatively good correspondence~\cite{quist2016,johansson2017}, and particle-replacement
models carry breakage into industrial compression crushers~\cite{cleary2015,delaney2015}, at a cost that one bonded-particle
study of a quarter section of an industrial cone put at $172$ and $404$ CPU hours for about four seconds of
operation~\cite{quist2012}. Automatic setting regulation and real-time
control raised the saleable yield of a crushing stage by $3.5\%$ (setting algorithm) and $4.2$ to $6.9\%$ (speed
algorithm) at full-scale aggregate plants~\cite{hulthen2010}, and the common opinion quoted by the same group is a yield
gain of $20$ to $25\%$ when liner wear is compensated~\cite{evertsson2015}; the optimised chamber geometry changes as the
liners wear~\cite{lindqvist2003}.

The practical gap that this report addresses lies between the first two families. The calibrations cited above are
judged on product size distributions, and some of them out of sample: Duarte et al.\ calibrate on two datasets and
validate on a third that the calibration did not use~\cite{duarte2021}, and Rocha et al.\ compare one global parameter
set, fitted to all ten surveys, with fits made survey by survey~\cite{rocha2024}. Neither reports the uncertainty of a prediction for a survey the
model has not seen, and neither holds the model to every survey in turn against baselines that use no physics; ten
surveys of one machine also leave little room for a free-form correction. A decision maker who reads a calibrated model at a new setting needs the error the model makes out of sample, a band that
covers the measured value at the stated rate, and a baseline that says how much of the skill is physics and how much
is the setting alone.

\subsection{What the product does}\label{sec:intro-product}

ChancaDEM is a static web studio with an offline pipeline behind it. A reader picks one of $27$ crushing cases (seventeen
circuit cases, fourteen across five machine classes and three controls, and the ten measured HP500 surveys of Rocha et
al.~\cite{rocha2024}), changes the closed-side setting, the throw, the speed, the feed and the ore, and reads the product
gradation, the throughput and the power from one engine: a Whiten population balance~\cite{whiten1972,napier1996} with an
Andersen and Napier-Munn classification~\cite{andersen1988}, a JKMRC $t_{10}$ energy-fineness relation with an Austin
appearance function~\cite{napier1996,austin1984}, a reduced form of the Evertsson capacity model~\cite{evertsson2000}, the
Bond power law~\cite{bond1952}, and, held out beside it, the Morrell size-specific crusher energy~\cite{morrell2009,gmsg}.
The engine runs in the browser in a TypeScript port held to the canonical Python engine by a shared reference of $\nrefpoints$
operating points. Two learned models, a multilayer perceptron that emulates the engine in the manner of Koh et
al.~\cite{koh2021} and a denoising autoencoder that scores how far a result lies from the training manifold, run in the
browser through ONNX Runtime. A projection of the hours until the product leaves a size target as the setting drifts
open, with the drift rate supplied by the user, completes the operating view.

\subsection{Contributions}\label{sec:intro-contrib}

The report documents every method of the product and the evaluation that holds them to measured data. Its contributions
are:
\begin{enumerate}
\item A complete account of a crusher engine (feed, breakage, classification, the Whiten solve, capacity, power, regimes
and an envelope check) with its equations, constants and numerical scheme, implemented twice and held equal to a largest
relative difference of $\refworst$ over $\nrefpoints$ reference points.
\item A strict leave-one-survey-out evaluation on the ten HP500 surveys of the nine predictors of Table~\ref{tab:bench} and of
the published throughput model of Rocha et al.\ in two variants, the same folds for every predictor, with a constant-mean baseline, a label-shuffle control, the optimism of a leaky evaluation measured beside the
strict one, and a nested $80\%$ predictive band whose coverage is computed out of sample, and a bootstrap over the surveys that gives every
held-out error and every difference between two predictors its interval.
\item The finding that a line on the closed-side setting alone has the lowest point estimate of the held-out power error on
these data ($12.78\%$, $95\%$ interval $\uqPLinelo$ to $\uqPLinehi$), not distinguishable on ten surveys from the free-form fit,
CrusherBayes and Morrell's energy with a no-load term, that the Bond draw of the calibrated engine carries little of the
measured power variation, and that the published throughput model of Rocha et al., run on the same surveys, is not
distinguishable from the calibrated engine on the feed rate.
\item CrusherBayes, a Bayesian calibration of the engine's eight fitted scalars with a structural discrepancy per target
and a differential-evolution sampler, whose held-out $80\%$ intervals cover eight and nine of the ten surveys (the power
coverage at the upper edge of the acceptable range) and are narrower than the residual band, and whose posterior shows which
scalars ten surveys contract and that the capacity base is contracted only through the prior scale of the discrepancy; the
method is reported in full in its own manuscript~\cite{crusherbayes2026} and summarised here.
\item An inventory of thirteen data contracts with their schemas, producers, consumers and missing-data behaviour, which
lets a reader replace any component and keep the rest.
\end{enumerate}

\subsection{Structure of the report}\label{sec:intro-structure}

Section~\ref{sec:problem} states the problem as an estimation problem with its inputs, outputs and assumptions, and
Section~\ref{sec:notation} collects the notation. Section~\ref{sec:data} describes the data sources and every data
contract. Section~\ref{sec:arch} describes the pipeline, the engine, the web and the dormant interface. Section~\ref{sec:methods}
gives one subsection per method, from the feed distribution to CrusherBayes. Section~\ref{sec:protocol} states the cases,
the metrics, the baselines and the acceptance criterion of each method; Section~\ref{sec:results} reports every result the
studio shows; Section~\ref{sec:validation} explains how the results are known to be correct. Sections~\ref{sec:discussion}
and~\ref{sec:limits} discuss what the results mean and what bounds them, Section~\ref{sec:repro} lists the commands that
rebuild every table and figure, and Section~\ref{sec:conclusions} concludes. The appendices hold the parameter tables, the derivation
of the integrated discrepancy of CrusherBayes, and the field tables of every contract schema. A reader interested only in the evaluation
on the measured surveys can read Sections~\ref{sec:m-hp500} to~\ref{sec:m-crushercal}, \ref{sec:m-crusherbayes} and~\ref{sec:results-bench}
to~\ref{sec:results-cb}; a reader who wants to reimplement the engine needs Sections~\ref{sec:m-feed} to~\ref{sec:m-bond} and the appendices.
